\section{Conclusion}
\label{sec:conclusion}
This study analyzed citation patterns in GSEs across political domains, revealing significant differences between GSE models and political parties in primary information source usage and highlighting the poisoning attack surface of GSEs.
Our method identified publisher attributes and quantified how publisher categories influenced answer generation by examining citation preferences and Citation Grounding Scores in GSE outputs.
Our results showed that primary information sources comprised $60\%$--$65\%$ in Japan versus $25\%$--$45\%$ in the U.S., and low-barrier sources comprised approximately $30\%$ of citations, indicating higher poisoning risk for U.S. political answers.


\section*{GenAI Usage Disclosure}
In this study, we used Claude Opus 4.7 to assist in producing visualization scripts and restructuring framework code to improve efficiency during the research stage, and to provide language refinement of author-written drafts to enhance clarity and conciseness during the writing stage. 
Every output was inspected and revised by the authors, who take full responsibility for the final content.